\scanpage{194}\mspage{50}
% archive sheet number: 194

et

\[
\begin{aligned}
&\text{Si }M'\hookrightarrow M^{\unsure{\deleted{\prime}|E|F}}\text{ est pl. fid, alors }M^{\unsure{\deleted{\prime}|E|\leftarrow}}\text{-connexe}\\
&\text{implique }M'\text{-connexe,}
\end{aligned}\tag*{(4.4.)}
\]

enfin, inversement

\[
\begin{aligned}
&\text{pour que }u\text{ soit }M\text{-connexe, il f. et s.}\\
&\text{que pour tte \uline{petite} sous-catégorie}\\
&\text{pleine }M'\text{ de }M,\ u\text{ soit }M'\text{ connexe.}
\end{aligned}\tag*{(4.5.)}
\]

Comme toute petite $M'$ se plonge dans $M'{}^\wedge\deleted{\unsure{,}}=\deleted{\unsure{M}}\mathrm{Ens}(M')$, il résulte de \unsure{.} (4.3)(4.4)(4.5) que \deleted{pour \unsure{tout}}

\[
\begin{aligned}
&\text{Si }u\text{ est Ens-connexe, il est }M\text{-connexe}\\
&\text{pour tte catégorie }M,
\end{aligned}\tag*{(4.6)}
\]

car c'est le cas si $M$ est petite par (4.3), donc pour tte $M$ par (4.5).

Disons qu'un foncteur $u$ \inserted{dans $\mathrm{Cat}$} est $M$-\uline{cofinal} \inserted{ou $M$-coasphérique} si \unsure{le|la} \deleted{\unsure{foncteur|fonction}} \unsure{hom|homo} fonctoriel

\[
\underset{\deleted{\unsure{*}}I}{\underset{\longleftarrow}{\lim}}\,F\longrightarrow
\underset{I'}{\underset{\longleftarrow}{\lim}}\,u^*(F)\tag*{(4.7)}
\]

(flèche dans $\underline{\mathrm{Hom}}(M(I),M)$ si les $\underset{\longleftarrow}{\lim}$ de types $I,I'$ existent dans $M$, sinon flèche dans $\underline{\mathrm{Hom}}(M(I),M^\wedge)$) est \deleted{\unsure{un}} iso. En se ramenant au cas de petites catégories $M$, \inserted{\unsure{chacune|disons|dûment}} plongées dans $M^\wedge$, on voit par le critère (4.2.bis) que

\[
u\text{ est }M\text{-connexe si pour tt }i\in\mathrm{Ob}\,I,\tag*{(4.8.)}
\]
